\section{Fixtures}

The L0 positive fixture is \texttt{l0-rename}. It is a one-event package whose
\texttt{LineageEventKind::Rename} must be one-to-one. The core test
\texttt{l0\_rename\_fixture\_verifies} checks that the lineage unit references
and lineage cardinality reports are present.

\texttt{l0-missing-unit} is the stale-unit negative fixture. It declares a
lineage event that references a unit not present in the relevant cycle context.
\texttt{rhist-core} rejects it with \texttt{LineageMissingUnit}, and
\texttt{rhist-io} surfaces the same core error when loading the directory.

\texttt{l1-split-merge} is the split/merge positive package. It contains two
lineage events and four crosswalk rows. The verifier checks both
\texttt{crosswalk\_unit\_refs} and \texttt{crosswalk\_weight\_sum}, proving
that the synthetic split and merge can be carried with exhaustive rational
weights.

\texttt{l1-bad-weights} is the bad-weight negative package. Its unit and
lineage references are shaped like the L1 case, but the exhaustive crosswalk
weights do not sum to one. The expected failure is
\texttt{RhistCoreError::CrosswalkWeightSum}.

\texttt{l2-three-cycle} is the three-cycle positive package. The tests assert
that it has three cycles, seven lineage events, and nine crosswalk rows, and
that its event kinds include rename, split, and merge. It is also the fixture
used by the RCOUNT RHIST-reference consumer test and by
\texttt{SYN\_RHIST\_L2\_PACKAGE\_HASH}.

\texttt{real-ri-tract-unchanged} is the real-source pressure fixture. It has
three cycles, two lineage events, and two crosswalk rows. The verifier checks
source reference resolution and \texttt{rhist-io} verifies that the preserved
source bytes match \path{sources/source-index.json}.

The hash tests cover the package substrate itself. One test confirms that
\texttt{package\_content\_hash} starts with \texttt{sha256:}, uses the RHIST
prefix, and ignores the declared hash field. Another changes a lineage
explanation and requires the hash to change; a third edits the manifest hash
and expects \texttt{PackageHashMismatch}.
